\section{Additional short equations: 73 operations at degree 54}
\label{sec:short-synthesis}

The incoming archives contain useful improvements that are absent from
the active P76 documentation: a 66-operation equation of degree 157,
a 72-operation equation of degree 80, and an ordinary gap equation with
75 operations at degree 44. We retain these as the comparison baseline.
The following synthesis adds the two tradeoffs
\[
 (\text{operations},\text{degree})=(73,54),\qquad (74,46),
\]
each using fifteen strictly positive integer variables. It combines the
incoming minimum-norm argument and squared-modulus reconstruction, and
regroups the main norm to save one further operation. The finite verifier
checks every literal arithmetic gate and exact leading homogeneous form.
The semantic compiler remains an identified inherited mathematical
dependency; no new proof-assistant or object-language ZFC certificate is
asserted in this supplement.

\subsection{The exact polynomial}

Use the complete package-31 coefficient tuple, with $B=2^{1064921}$,
$\beta=B-1$, and its already shifted $M_F$. The tuple and all four binary
payloads are included in \texttt{short/coefficients31/}.
The fifteen positive unknowns are
\[
 J,F,\alpha,z,f,h,i,T,s,w,\tau,\eta,\zeta,e,\gamma.
\]
Define polynomial abbreviations, with no additional unknowns:
\begin{align*}
 q&=\beta J+1,& C&=q-F-\alpha,\\
 X&=wq,&Y&=sq^3,&E&=XY,\\
 k&=\eta+\zeta,&c&=kY+\eta,&a&=E+Y,\\
 H&=4a+3,&\Delta&=a^2+H,&D&=X+ac+\gamma H,\\
 L&=E(kY),&
 R&=((q-1)(q-F)+\alpha)(q^2-1)+(M_C+qM_F)J,\\
 V&=f^2(cT-R)-c,&U&=\Delta(f^2-1),&I&=ic^2.
\end{align*}
The five factors and strong residual are
\begin{align*}
 N_0&=\tau^2-L(L+k),&N_1&=D^2-\Delta c^2,\\
 N_a&=(V+e)^2-Ue(2V+e),&W_s&=I^2-U,\\
 N_k&=k-hE-R,&N_t&=(K_0+w)C+(q-F)-z(q-1).
\end{align*}
The main short equation is
\begin{equation}\label{eq:short54}
 \boxed{Q_{54}=N_0^2+N_a^2+W_s^2+(N_1N_kN_t-1)^2-2=0.}
\end{equation}
The lower-degree alternative is
\begin{equation}\label{eq:short46}
 \boxed{Q_{46}=N_0^2+N_1^2+N_a^2+W_s^2+(N_kN_t-1)^2-3=0.}
\end{equation}
The subscripts denote exact degree. Their literal evaluators have 73
and 74 operations, respectively. The principal evaluator has
39 multiplications, 19 additions and 15 subtractions; the alternative
has one more addition. Multiplication by a fixed numeral is charged.

\subsection{The minimum argument and complete local equivalence}

The following facts hold on every positive assignment, before imposing
either equation:
\begin{equation}\label{eq:min-blocks}
 N_0^2\geq1,\qquad N_a^2+W_s^2\geq1.
\end{equation}
Equality cases are $N_0=1$ and $(N_a,W_s)=(1,0)$, respectively.
For the first assertion, write $n=XY^2\geq2$. The discriminant $n(n+1)$
is nonsquare, so $N_0\ne0$. A hypothetical solution of
$x^2-n(n+1)t^2=-1$ has
$nt<|x|<(n+\tfrac12)t$. Multiplication by
$2n+1-2\sqrt{n(n+1)}$ produces a new positive second coordinate
$(2n+1)t-2|x|<t$, contradicting minimality. Thus $N_0\ne-1$.

For the second assertion put $y_0=V+e$. The exact identity
\[
 N_a=UV^2-(U-1)y_0^2
\]
is immediate. If $U=0$, then $W_s=I^2\geq81$. Otherwise $U\geq2$,
and $N_a=0$ would force the consecutive integers $U,U-1$ to be
squares, unless $V=y_0=0$. The latter contradicts $e>0$. Hence
$N_a\ne0$. Equality in the second block forces $W_s=0$ and
$N_a=\pm1$. Now $U=I^2$, so modulo four $N_a$ is congruent to
$y_0^2$ when $I$ is even and to $V^2$ when $I$ is odd. Thus $N_a=-1$
is impossible. This proves all assertions in~\eqref{eq:min-blocks}.

At a zero of~\eqref{eq:short54}, both minimum blocks attain their
minima and the remaining square vanishes. Hence
\[
 N_0=N_a=1,\qquad W_s=0,\qquad N_1N_kN_t=1.
\]
The last equation makes $N_1$ an integer unit. Since
$\Delta=(a+2)^2-1\equiv0$ or $3\pmod4$, square residues exclude
$N_1=-1$. Thus $N_1=1$ and $N_kN_t=1$. These are precisely the
equality conditions of~\eqref{eq:short46}; conversely they imply both
equations. This is a complete tuplewise equivalence, uniform in the
fixed ports with $\beta>0$.

The saving is transparent at the circuit level. Starting from $Q_{46}$,
remove the $N_1^2$ gate and its summation gate, insert one multiplication
of $N_1$ into the unit product, and change the final constant from three
to two. This saves exactly one binary operation.

\subsection{Squared-modulus completeness for the ordinary strong form}

The historical gap construction used
$V_{\mathrm{old}}=c(T_{\mathrm{old}}f-1)-Rf^2$.
The identity
\[
 V=f^2(cT-R)-c=V_{\mathrm{old}}\big|_{T_{\mathrm{old}}=fT}
\]
proves soundness immediately. Completeness requires reconstructing an
auxiliary Pell ordinate, rather than assuming divisibility of every
old quotient.

A halting computation supplies a canonical normalized witness.
Transfer it to ordinary coordinates by
$i_{\mathrm{ordinary}}=\Delta i_{\mathrm{normalized}}$.
The inherited decoding supplies $c=\psi_A(R)$,
$A=a+2$, $R\equiv3\pmod4$, and
\[
 I^2=\Delta(f^2-1),\qquad c\mid I,
 \qquad f^2\equiv1\pmod c,
 \qquad \gcd(c,f)=1.
\]
For $R=2r+1$ the Pell polynomials satisfy
\[
 \chi_t(R)=tQ_r(t^2),\quad
 Q_r(0)=(-1)^rR,\quad
 Q_r(1-A^2)=(-1)^r\psi_A(R).
\]
These follow by the elementary second-order Pell recurrences. Since $r$
is odd, take
\[
 V=\frac{\chi_I(R)}{I},\qquad y=\psi_I(R).
\]
Both are positive integers, $V>1$, and their norm is one. The displayed
polynomial identities give $V\equiv-c\pmod{f^2}$ and
$V\equiv-R\pmod c$. Consequently
\[
 T=\frac{V+c+Rf^2}{cf^2}>0
\]
is an integer and satisfies the required quotient formula. Finally
\[
 (I^2-1)(y^2-V^2)=V^2-1>0
\]
implies $e=y-V>0$. The other coordinates remain unchanged. This proves
existence equivalence to the authenticated ordinary source and hence,
using its inherited compiler and K215 calibration, the ordinary
mathematical implications formalizable in ZFC:
\begin{align*}
 \text{$Q_{54}$ has no positive zero}&\longleftrightarrow
 \operatorname{Con}(\mathrm{ZFC}),\\
 \text{$Q_{46}$ has no positive zero}&\longleftrightarrow
 \operatorname{Con}(\mathrm{ZFC}).
\end{align*}

\subsection{Exact degrees and bounded comparison}

The factor degrees are $22,18,23,22,7,2$ for
$N_0,N_1,N_a,W_s,N_k,N_t$, respectively. The main-norm cancellation is
the identity, with $Z=X+\gamma H$,
\[
 (ac+Z)^2-(a^2+H)c^2=Z^2+2acZ-Hc^2.
\]
Writing $Q=\beta J$, $k_0=\eta+\zeta$ and
$C_1=Q-F-\alpha$, the highest homogeneous forms are
\begin{align*}
 [Q_{54}]_{54}
 &=64\gamma^2h^2w^6s^8k_0^2Q^{30}(wC_1-zQ)^2,\\
 [Q_{46}]_{46}
 &=4w^4s^6k_0^2T^2f^8e^2Q^{22}.
\end{align*}
Their respective coefficients of
$J^{32}\gamma^2h^2w^6s^8\eta^2z^2$ and
$J^{22}w^4s^6\eta^2T^2f^8e^2$ are
$64\beta^{32}$ and $4\beta^{22}$, strictly positive for the supplied
tuple. The exact symbolic audit verifies these full leading forms,
not merely a numerical specialization.

These points improve the local degree/operation frontier between the
reviewed $72/80$ and $75/44$ candidates. They do not improve the existing
66-operation record, the degree-44 point, or the smallest-variable
constructions. The proof files and source receipts make this a bounded
repository comparison, without a claim of worldwide priority or global
minimality.
